\documentclass[11pt]{article}

\usepackage[a4paper,margin=1in]{geometry}
\usepackage{amsmath,amssymb,amsthm}
\usepackage{graphicx}
\usepackage{booktabs}
\usepackage{xcolor}
\usepackage{verbatim}
\usepackage[colorlinks=true,linkcolor=blue!55!black,citecolor=blue!55!black,
            urlcolor=blue!55!black]{hyperref}

\definecolor{codecm}{RGB}{120,130,140}

% render all verbatim code blocks in a smaller monospace font
\makeatletter
\renewcommand{\verbatim@font}{\normalfont\ttfamily\small}
\makeatother

\newtheorem{theorem}{Theorem}
\newtheorem{proposition}{Proposition}
\newtheorem{definition}{Definition}
\newtheorem{remark}{Remark}

\newcommand{\R}{\mathbb{R}}
\newcommand{\E}{\mathbb{E}}
\newcommand{\Tcheb}{T}
\newcommand{\code}[1]{\texttt{#1}}
\newcommand{\half}{\tfrac{1}{2}}

\title{\vspace{-1.2cm}\bfseries Chebyshev Collocation, Smolyak Sparse Grids,
and Two-Branch Quadrature\\[4pt]
\large A Textbook Companion to the Global Solver in \code{code/bocola2016/}}
\author{Computational methods note --- Bocola (2016) replication package}
\date{\today}

\begin{document}
\maketitle

\begin{abstract}
\noindent This note explains, at a graduate textbook level, the numerical
machinery that the \code{code/bocola2016/} package uses to solve a recursive
competitive equilibrium with priced sovereign default risk. Three ideas do
the work. \textbf{Chebyshev polynomial approximation} turns an unknown
policy function into a short vector of coefficients with near-optimal,
spectrally accurate interpolation on carefully chosen nodes.
\textbf{Smolyak sparse grids} extend that construction to six state
variables without paying the tensor-product ``curse of dimensionality.''
And a \textbf{two-branch Gauss--Hermite quadrature} evaluates the
conditional expectations inside the agents' Euler equations, splitting the
future into a no-default and a default continuation weighted by a priced
probability. For each idea we give the underlying theory, the specific
numerical hazards it carries, and the exact place in the code where it
lives. The recurring theme is that every one of these tools is a controlled
approximation whose error is well understood only inside a region the model
actually visits---and that respecting that region (clipping, rotation,
smoothing, best-iterate tracking) is what separates a working global solver
from a diverging one.
\end{abstract}

\tableofcontents

\vspace{0.5cm}
\section{The problem being solved}
\label{sec:problem}

The model is a detrended recursive competitive equilibrium. The aggregate
state is the six-vector
\[
S \;=\; (\,\underbrace{K,\;B,\;P}_{\text{endogenous}},\;
\underbrace{\Delta z,\;g,\;s}_{\text{exogenous}}\,)\in\R^{6},
\]
capital $K$, the sovereign bond stock $B$, banks' promised deposit payment
$P$, the technology growth innovation $\Delta z$, the fiscal shock $g$, and
the \emph{risk factor} $s$ that governs the priced probability of a
sovereign default next period. A solution is not a number or a finite path;
it is a set of \emph{decision rules}---functions of the state---for
consumption $C(S)$, the deposit rate $R(S)$, the bond price $Q_B(S)$, and
the banker's marginal value of wealth $\alpha(S)$ (written \code{av} in the
code). These functions must satisfy, at \emph{every} point of the state
space, a system of functional equations: three Euler equations (deposits,
capital, sovereign bonds), the recursion for $\alpha(S)$, and an
occasionally-binding leverage constraint.

Two facts force a \emph{global, nonlinear} method rather than a local
(perturbation) one. First, the leverage constraint binds only
occasionally---it is slack in normal times and tightens sharply in
stress---so the policy functions have a \emph{kink} that a Taylor expansion
around the balanced-growth path cannot capture. Second, the object of
interest is precisely the behaviour far from that path (a sovereign-risk
crisis). A global method represents each unknown function by a finite basis
expansion and pins down the coefficients by forcing the residual of the
functional equations to zero at a finite set of \emph{collocation nodes}.
This is the \emph{projection} (or \emph{weighted-residual}) approach. Three
questions organize the rest of the note:

\begin{enumerate}
\item \textbf{Which basis and which nodes?} $\;\to\;$ Chebyshev polynomials
      on Chebyshev--Lobatto nodes (\S\ref{sec:cheb}).
\item \textbf{How to do this in six dimensions?} $\;\to\;$ Smolyak sparse
      grids (\S\ref{sec:smolyak}).
\item \textbf{How to evaluate the expectations that make the equations
      \emph{Euler} equations?} $\;\to\;$ two-branch Gauss--Hermite quadrature
      (\S\ref{sec:branch}).
\end{enumerate}

\noindent Section~\ref{sec:risks} is devoted entirely to the numerical
\emph{risks}---the ways each tool fails, and how the code defends against
them.


\section{Univariate Chebyshev approximation}
\label{sec:cheb}

Everything multidimensional in this code is built from one univariate
object, so we start there. The goal is to approximate a smooth function
$f:[-1,1]\to\R$ by a polynomial $p_n$ of degree $n$ such that (i) the error
is as small as possible, (ii) the map from function values to coefficients
is well conditioned, and (iii) both are cheap to evaluate.

\subsection{Chebyshev polynomials of the first kind}

\begin{definition}
The Chebyshev polynomials of the first kind are defined on $[-1,1]$ by
\[
\Tcheb_k(x) \;=\; \cos\!\big(k\arccos x\big), \qquad k=0,1,2,\dots
\]
\end{definition}

\noindent The definition looks trigonometric but $\Tcheb_k$ is a genuine
polynomial of degree $k$: with $x=\cos\theta$, $\cos(k\theta)$ expands into
powers of $\cos\theta$. Equivalently, and this is exactly how the code
computes them, they obey the three-term recurrence
\begin{equation}
\Tcheb_0(x)=1,\qquad \Tcheb_1(x)=x,\qquad
\Tcheb_{k}(x)=2x\,\Tcheb_{k-1}(x)-\Tcheb_{k-2}(x).
\label{eq:recur}
\end{equation}
This recurrence is the entire content of \code{chebyshev\_basis\_1d} in
\code{smolyak.py}:
\begin{verbatim}
def chebyshev_basis_1d(x, max_deg):
    T = np.empty((x.size, max_deg + 1))
    T[:, 0] = 1.0
    if max_deg >= 1:
        T[:, 1] = x
    for k in range(2, max_deg + 1):
        T[:, k] = 2.0 * x * T[:, k - 1] - T[:, k - 2]
    return T
\end{verbatim}
The recurrence is used rather than the monomial form $\sum a_k x^k$ because
the monomial (Vandermonde) representation is catastrophically
ill-conditioned---the basis functions $x^k$ become nearly linearly
dependent as $k$ grows---whereas \eqref{eq:recur} is numerically stable.

\begin{proposition}[Orthogonality]
Under the weight $w(x)=(1-x^2)^{-1/2}$,
\[
\int_{-1}^{1}\frac{\Tcheb_j(x)\,\Tcheb_k(x)}{\sqrt{1-x^2}}\,dx
=\begin{cases}0 & j\neq k,\\ \pi & j=k=0,\\ \pi/2 & j=k\ge 1.\end{cases}
\]
\end{proposition}
This is the change of variables $x=\cos\theta$ applied to the elementary
orthogonality of $\{\cos k\theta\}$ on $[0,\pi]$. Orthogonality is what makes
a Chebyshev expansion a \emph{cosine series in disguise}: approximating
$f(x)$ by $\sum_k c_k \Tcheb_k(x)$ is the same as approximating the even,
$2\pi$-periodic function $f(\cos\theta)$ by a Fourier cosine series. All the
good behaviour of Fourier series on periodic functions transfers, without
their bad behaviour at non-periodic endpoints.

\subsection{The nodes: why not equispaced?}
\label{sec:runge}

A degree-$n$ interpolant is pinned by $n+1$ nodes. The naive choice---evenly
spaced points---is a disaster for polynomial interpolation. The classic
counterexample is Runge's function $f(x)=1/(1+25x^2)$: as the degree rises,
the equispaced interpolant develops ever-larger oscillations near
$x=\pm1$ and \emph{diverges} in the sup norm, even though $f$ is perfectly
smooth (Figure~\ref{fig:runge}, left). This is the \emph{Runge phenomenon},
and it reflects a deep fact: interpolation on equispaced nodes is
exponentially ill-conditioned.\footnote{See the discussion of
Runge/Lebesgue in Trefethen's \emph{Approximation Theory and Approximation
Practice}, and the survey references collected at the end of this note.}

\begin{figure}[h]
\centering
\includegraphics[width=\textwidth]{runge.pdf}
\caption{Degree-14 polynomial interpolation of Runge's function. Equispaced
nodes (left) oscillate wildly near the endpoints; Chebyshev--Lobatto nodes
(right) track the function. Same degree, same function---only the node
placement differs.}
\label{fig:runge}
\end{figure}

The cure is to cluster nodes toward the endpoints. The \emph{Chebyshev
extrema}, also called \emph{Chebyshev--Lobatto} or \emph{Clenshaw--Curtis}
points, do exactly this:
\begin{equation}
x_k \;=\; -\cos\!\Big(\frac{\pi k}{m-1}\Big),\qquad k=0,\dots,m-1,
\label{eq:extrema}
\end{equation}
which are the projections onto the $x$-axis of $m$ equally spaced points on
the upper unit semicircle (Figure~\ref{fig:nodes}). This is the
\code{\_level\_points} routine in \code{smolyak.py}. Their density near
$\pm1$ scales like $(1-x^2)^{-1/2}$---exactly the Chebyshev weight---so the
interpolation problem becomes well conditioned.

\begin{figure}[h]
\centering
\includegraphics[width=0.72\textwidth]{nodes.pdf}
\caption{The Chebyshev--Lobatto nodes \eqref{eq:extrema} are equally spaced
\emph{angles} on the semicircle projected down onto $[-1,1]$; the projection
clusters them near the endpoints, which is what tames the Runge
phenomenon.}
\label{fig:nodes}
\end{figure}

\subsection{The Lebesgue constant: near-optimality}
\label{sec:lebesgue}

The quality of a set of nodes is measured by the \emph{Lebesgue constant}
$\Lambda_n$, the operator norm of the interpolation map from
$C[-1,1]$ to itself. It bounds how far interpolation can be from the best
possible polynomial approximation $p_n^\star$:
\begin{equation}
\|f-p_n\|_\infty \;\le\; \big(1+\Lambda_n\big)\,\|f-p_n^\star\|_\infty .
\label{eq:lebesgue}
\end{equation}
A small $\Lambda_n$ means ``interpolating is almost as good as doing the best
you possibly could.'' The contrast is stark:
\[
\Lambda_n^{\text{equispaced}} \sim \frac{2^{n+1}}{e\,n\log n}
\quad(\text{exponential}),
\qquad
\Lambda_n^{\text{Chebyshev}} \sim \frac{2}{\pi}\log n
\quad(\text{logarithmic}).
\]
Chebyshev nodes are not \emph{exactly} optimal---the minimizing nodes have
no closed form---but they capture essentially all of the available accuracy,
with a Lebesgue constant that grows so slowly ($\log n$) that it is
irrelevant in practice. This is the precise sense of the phrase ``Chebyshev
interpolation is near-minimax.''

\subsection{Convergence rate: spectral accuracy}
\label{sec:spectral}

How fast does the error fall as we add degrees? The answer depends entirely
on the smoothness of $f$:

\begin{theorem}[Rate governed by smoothness]
Let $p_n$ be the degree-$n$ Chebyshev interpolant of $f$ on $[-1,1]$.
\begin{itemize}
\item If $f$ has $\nu$ continuous derivatives with bounded variation, then
      $\|f-p_n\|_\infty = O(n^{-\nu})$ (\emph{algebraic}).
\item If $f$ is analytic in a Bernstein ellipse with parameter $\rho>1$
      around $[-1,1]$, then $\|f-p_n\|_\infty = O(\rho^{-n})$
      (\emph{geometric / spectral}).
\end{itemize}
\end{theorem}

\noindent Figure~\ref{fig:conv} shows both regimes on Chebyshev nodes: the
analytic target $e^{\sin 3x}$ hits machine precision by degree $\sim30$
(straight line on a log scale = geometric decay), while the kinked target
$|x-0.2|$ crawls down algebraically. This dichotomy is the single most
important fact for the model at hand. Where its policy functions are smooth,
a handful of Chebyshev degrees delivers near-machine accuracy. But the
occasionally-binding leverage constraint injects a kink into the true
policy, and \emph{at that kink the spectral rate collapses to algebraic}---a
polynomial basis chasing a non-smooth target produces Gibbs-type ringing.
This is not a coding defect; it is a theorem about polynomials, and it
motivates the smoothing devices of \S\ref{sec:risks}.

\begin{figure}[h]
\centering
\includegraphics[width=0.62\textwidth]{convergence.pdf}
\caption{Maximum interpolation error versus polynomial degree on Chebyshev
nodes. Analytic functions converge geometrically (spectral accuracy); a
kink forces algebraic convergence. The model's occasionally-binding
constraint puts a kink in the policy, so the second curve is the relevant
warning.}
\label{fig:conv}
\end{figure}

\subsection{Collocation as a linear solve}
\label{sec:collocation}

Fix the degrees $\{0,1,\dots,n\}$ and the nodes $\{x_k\}_{k=0}^{n}$. Writing
the interpolant as $p(x)=\sum_{j} c_j \Tcheb_j(x)$ and demanding
$p(x_k)=f(x_k)$ at every node gives the square linear system
\begin{equation}
\underbrace{\big[\Tcheb_j(x_k)\big]_{k,j}}_{\text{basis matrix }\Phi}
\; c \;=\; f,
\label{eq:collocation}
\end{equation}
so the coefficients are $c=\Phi^{-1}f$. Because $\Phi$ depends only on the
grid, not on $f$, the code factorizes it \emph{once} (an LU factorization)
and reuses it every time it needs coefficients from values. In
\code{SmolyakGrid} this is the pair \code{fit} (solve for $c$) and
\code{eval} (form $\Phi(x)c$ at arbitrary points), with the LU factor
\code{self.\_lu} built in the constructor. Two consequences worth
internalizing:

\begin{itemize}
\item \textbf{Interpolation is exact at the nodes by construction.} The
      interpolant reproduces $f$ at every collocation point to rounding
      error; all approximation error lives \emph{between} nodes.
\item \textbf{Evaluating the interpolant outside $[-1,1]$ is extrapolation},
      and Chebyshev extrapolation grows like $|\Tcheb_n(x)|\sim
      |x+\sqrt{x^2-1}|^{\,n}$---explosively. This is the single most
      dangerous property of the whole scheme and is treated in
      \S\ref{sec:risk-extrap}.
\end{itemize}

Finally, the code maps the model's natural state box $[\ell,h]$ onto the
reference interval by the affine transform
$u=2(x-\ell)/(h-\ell)-1$ (\code{to\_unit}/\code{from\_unit}); all Chebyshev
machinery lives in the $u\in[-1,1]$ coordinates.


\section{Smolyak sparse grids}
\label{sec:smolyak}

Section \ref{sec:cheb} solved the one-dimensional problem. The model has
six state dimensions. The obvious extension---a \emph{tensor product} of
one-dimensional grids---is computationally fatal, and the fix is Smolyak's
1963 construction, brought into economics by Krueger and Kubler (2004) and
refined by Judd, Maliar, Maliar and Valero (2014).

\subsection{The tensor-product curse}

A tensor grid with $m$ points per dimension has $m^{d}$ points in $d$
dimensions. With a modest $m=5$ and the model's $d=6$ that is
$5^6=15{,}625$ collocation nodes---and every node requires solving the
agents' nonlinear system and a full expectation, at every outer iteration.
At $m=9$ it is over half a million. The count grows \emph{geometrically in
the dimension}: this is Bellman's \emph{curse of dimensionality}
(Figure~\ref{fig:growth}, grey line). Tensor grids are simply not an option
past three or four states.

\begin{figure}[h]
\centering
\includegraphics[width=0.6\textwidth]{growth.pdf}
\caption{Node count versus state dimension. A tensor product explodes
geometrically; the Smolyak grids at levels $\mu=2,3$ grow polynomially.
At the model's $d=6$ the Smolyak grid has $85$ ($\mu{=}2$) or $389$
($\mu{=}3$) nodes---versus tens of thousands for a comparable tensor grid.}
\label{fig:growth}
\end{figure}

\subsection{Nested one-dimensional levels}

The Smolyak idea rests on a sequence of \emph{nested} one-dimensional grids.
Define level $i$ to have $m_i$ Chebyshev--Lobatto points, with
\begin{equation}
m_1=1,\qquad m_i = 2^{\,i-1}+1 \ \ (i\ge2)\ \Rightarrow\ m=1,3,5,9,17,\dots
\label{eq:levels}
\end{equation}
The doubling rule \eqref{eq:levels} makes the grids \emph{nested}: every
point of level $i$ is also a point of level $i{+}1$
(\code{\_level\_points}). Nestedness is what lets Smolyak add resolution
without recomputing anything, and it is why Clenshaw--Curtis/Chebyshev
extrema (rather than, say, Gauss points) are the natural building block. The
code works with the \emph{incremental} sets---the points and the polynomial
degrees that are \emph{new} at level $i$:
\begin{verbatim}
def _new_points(i):        # points introduced at level i (disjoint across i)
    if i == 1: return np.array([0.0])
    if i == 2: return np.array([-1.0, 1.0])
    return _level_points(i)[1::2]     # odd positions absent from level i-1

def _new_degrees(i):       # Chebyshev degrees introduced at level i
    if i == 1: return np.array([0])
    if i == 2: return np.array([1, 2])
    return np.arange(2**(i-2)+1, 2**(i-1)+1)
\end{verbatim}
Crucially $|\,$new points$\,|=|\,$new degrees$\,|$ at every level. This is
the bookkeeping that keeps the eventual multidimensional basis matrix
\emph{square}.

\subsection{The Smolyak combination rule}
\label{sec:combination}

In $d$ dimensions, index a tensor block by a multi-index
$\mathbf{i}=(i_1,\dots,i_d)$ of per-dimension levels. A full tensor grid at
level $\mu{+}1$ would take \emph{all} blocks with each $i_j\le\mu{+}1$.
Smolyak instead keeps only the blocks whose \emph{total} level is bounded:
\begin{equation}
\text{keep } \mathbf{i}
\quad\Longleftrightarrow\quad
\sum_{j=1}^{d}\big(i_j-1\big)\;\le\;\mu ,
\label{eq:smolyak}
\end{equation}
where $\mu\ge0$ is the \emph{approximation level}. The grid is the union of
the tensor products of the \emph{new-point} sets over all kept $\mathbf{i}$;
the basis is the union of the tensor products of the \emph{new-degree} sets
over the same $\mathbf{i}$. This is precisely \code{\_multi\_indices} +
the assembly loop in \code{SmolyakGrid.\_\_init\_\_}:
\begin{verbatim}
for i_vec in _multi_indices(d, mu, mu_vec):
    axes_p = [_new_points(i)  for i in i_vec]      # nested points
    axes_d = [_new_degrees(i) for i in i_vec]      # matching degrees
    pts.append(cartesian_product(axes_p))          # tensor of the *new* sets
    degs.append(cartesian_product(axes_d))
\end{verbatim}
Because points and degrees are added in lock-step, the assembled basis
matrix $\Phi=\big[\prod_j \Tcheb_{\text{deg}_j}(u_j)\big]$ evaluated at the
grid points is \emph{square and invertible}: the interpolant is \emph{exact
at every Smolyak node}, and \code{fit}/\code{eval} are again a single
LU solve. Selecting only ``low total level'' blocks is exactly what drops
the expensive high-resolution-in-many-dimensions-at-once corners that a
tensor grid wastes its budget on. In two dimensions the effect is visible
directly (Figure~\ref{fig:2d}): the Smolyak grid is a sparse cross-shaped
subset of the full tensor grid.

\begin{figure}[h]
\centering
\includegraphics[width=\textwidth]{smolyak2d.pdf}
\caption{In two dimensions: the full $9\times9$ tensor grid (left) versus
the Smolyak grids at $\mu=2$ ($13$ nodes) and $\mu=3$ ($29$ nodes). The
sparse grids keep the axes and low-order cross terms and discard the
dense interior---the source of the savings.}
\label{fig:2d}
\end{figure}

\subsection{What Smolyak approximates, and anisotropy}

The resulting interpolant is exact for the \emph{hyperbolic-cross} polynomial
space---all monomials $\prod_j x_j^{a_j}$ with a bounded weighted total
degree---rather than the full tensor space. Practically, $\mu=2$ captures
all quadratic cross terms (a complete second-order approximation plus
selected higher terms); $\mu=3$ adds cubic interactions. The package's gate
tests confirm exactly this: on-grid exactness to $10^{-10}$, complete
quadratic exactness including cross terms, and geometric error decay
$10^{-1}\!\to\!10^{-2}\!\to\!5\times10^{-4}$ as $\mu$ runs $2,3,4$ on a
smooth 6-D test function.

The code also supports \emph{anisotropic} levels through \code{mu\_vec}: a
per-dimension cap $i_j-1\le \mu_j$ in \code{\_multi\_indices}. This is the
Judd--Maliar--Valero refinement---spend resolution on the dimensions where
the policy actually curves and economize on the rest. The model does not
lean on anisotropy heavily, but the hook is there.

\subsection{The PCA-rotated grid}
\label{sec:rotated}

One more construction deserves mention because it is where sparse-grid
practice meets the specific economics. A box aligned with the physical axes
$(K,B,P)$ badly overcovers the state space: the ergodic cloud is elongated
along a diagonal (more capital $\Rightarrow$ more deposits $\Rightarrow$
higher obligations), so an axis-aligned box has deeply \emph{infeasible}
anti-correlated corners (low assets, high obligations) that the model never
visits but that the collocation fit is nonetheless forced to represent.
\code{RotatedGrid} (\code{rotated\_grid.py}) diagonalizes the ergodic
covariance of $(K,B,P)$, builds the Smolyak box in the principal-component
coordinates, and presents a physical-coordinate interface so the rest of the
code is unchanged. Rotating cut the infeasible-corner fraction from
$\sim48\%$ to $\sim16\%$ and delivered machine precision on-grid. It is a
concrete instance of the general lesson in \S\ref{sec:risks}: a global
polynomial basis is only as good as the region you ask it to cover.


\section{Numerical risks in the computation}
\label{sec:risks}

Chebyshev--Smolyak collocation is powerful but not forgiving. This section
catalogues the ways it goes wrong and the specific defenses in the code. A
PhD reader should treat this as the ``what will bite you'' section.

\subsection{Extrapolation blow-up}
\label{sec:risk-extrap}

The gravest hazard. Outside the box, Chebyshev polynomials grow like
$\rho^{\,n}$; a rule fitted on $[-1,1]$ and evaluated at $u=1.2$ can return
a \emph{negative} consumption or a nonsensical bond price, which then feeds
back through the expectation and produces \code{NaN}s. The danger is acute
here because the expectation step evaluates next-period rules at
\emph{stochastically generated} future states that can wander outside the
current box. The defense is to \emph{clip} the evaluation point back into
the box before every interpolation:
\begin{verbatim}
S_eval = grid.clip(S_next)          # never extrapolate the Chebyshev rule
Cn  = np.maximum(grid.eval(coef["C"],  S_eval), 1e-8)
QBn = np.maximum(grid.eval(coef["QB"], S_eval), 1e-4)
\end{verbatim}
Clipping is a \emph{documented, deliberate} projection (\code{clip} is a
public method, not a silent guard), and the numeric floors
($10^{-8}$, $10^{-4}$) are a second line of defense against the residual
overshoot near the boundary. The cost is a small bias when the true state
genuinely leaves the box; the package controls it by choosing the box wide
enough that the ergodic path is $99.3\%$ interior, so clipping essentially
never binds on a simulated path.

\subsection{Conditioning and the sparse-grid Lebesgue constant}

The one-dimensional Lebesgue constant is a gentle $\log n$
(\S\ref{sec:lebesgue}), but on Smolyak grids it grows faster---roughly
$(\log n)^{d}$-like behaviour in $d$ dimensions---so the amplification of
interpolation error, and the conditioning of the basis matrix $\Phi$,
degrade as the level $\mu$ or the dimension rises. In practice this is why
the code lives at $\mu=2$ (occasionally $\mu=3$): high levels buy formal
accuracy but a worse-conditioned solve. The stable Chebyshev
\emph{recurrence} (not the monomial basis) and the once-and-for-all LU
factorization keep the conditioning as good as the node set allows, but they
cannot repair an intrinsically hard node set---hence the preference for
modest levels plus a well-placed box over brute-force refinement.

\subsection{Aliasing and non-interpolatory error}

A sparse grid is exact only on the hyperbolic-cross polynomial space. Any
content of the true policy outside that space is \emph{aliased}---
misattributed to the retained basis functions. For a smooth policy this is
negligible (the aliased content is tiny), but two situations amplify it:
(i) high-frequency features from a near-kink, and (ii) using the same sparse
construction for \emph{integration}, where the Smolyak/Clenshaw--Curtis
weights can be \emph{negative}, so a sparse quadrature of a sharply peaked
integrand can return a value outside the range of the integrand. The code
sidesteps (ii) entirely by \emph{not} using the sparse grid to integrate:
expectations are done with a dedicated Gauss--Hermite tensor rule
(\S\ref{sec:branch}), whose weights are strictly positive.

\subsection{The residual curse of dimensionality}

Smolyak defeats the \emph{geometric} curse but not all cost growth. The node
count still rises polynomially in $d$ (Figure~\ref{fig:growth}), and
Krueger--Kubler and Malin--Krueger--Kubler document that beyond roughly
twenty state variables even sparse grids become expensive. Six states is
comfortably inside the sweet spot; the point for the reader is that
``Smolyak breaks the curse'' is a statement about the \emph{exponent}, not a
claim of free lunch.

\subsection{Non-smoothness: kinks, Gibbs, and Fischer--Burmeister}
\label{sec:risk-kink}

This is the model-specific hazard. The occasionally-binding leverage
constraint,
\[
\mu \;=\; \max\!\Big\{\,1-\frac{\E[\hat\Lambda']\,R\,N}{\lambda\,(Q_KK'+Q_BB')}\,,\;0\Big\},
\]
puts a $\max$-kink into the policy. Two distinct problems follow.
First, as a matter of \emph{approximation theory}, a global polynomial
basis converges only algebraically at a kink and produces Gibbs oscillations
nearby (\S\ref{sec:spectral})---the sparse grid \emph{cannot} represent the
constraint boundary sharply. Second, as a matter of \emph{solving}, the
$\max$ makes the finite-difference Jacobian of the outer nonlinear solve
noisy and stalls Newton. The code's answer, imported from the two-country
solver, is \emph{Fischer--Burmeister} complementarity smoothing: make $\mu$
an explicit unknown and replace the $\max$ with the equation
\begin{equation}
\varphi(\mu,\text{slack}) \;=\; \mu+\text{slack}-\sqrt{\mu^2+\text{slack}^2}\;=\;0,
\label{eq:fb}
\end{equation}
which is zero exactly when $\mu\ge0$, $\text{slack}\ge0$, and
$\mu\cdot\text{slack}=0$---the complementarity conditions---and is smooth
everywhere except the single point $(0,0)$ (\code{\_fb} in
\code{newton\_collocation.py}). A companion trick smooths a second kink
\emph{inside} the expectation: the hard investment floor
$\mathbf{1}\{I'>0\}$ is replaced by a softplus that equals $I'$ for feasible
investment to $10^{-12}$ and bends smoothly to a floor only near $I'=0$, so
the Jacobian stays well-defined at the quadrature nodes. The general
principle---any new $\max/\min/\mathbf{1}\{\cdot\}$ on a Newton-solved
quantity needs FB-style smoothing or the finite-difference Jacobian
stalls---is a recurring rule in this codebase.

\subsection{Boundary infeasibility and near-unit-root states}

The model's capital has a near-unit-root ($\sim0.98$) dynamic, so the
ergodic set is \emph{wide}. A box wide enough to contain it necessarily
includes corner states that are economically infeasible (negative net
worth, $\mu\to1$, the banker value function undefined). There, the
pointwise nonlinear solve has no sensible solution, and if those garbage
values enter the global fit they poison the interpolant \emph{everywhere}
through the coupling of the collocation system. The code deploys three
countermeasures, each addressing one facet:
\begin{itemize}
\item \textbf{Anchoring} (\code{make\_residual}, \code{feasible\_mask}):
      infeasible corners are held at their warm-start values via an anchor
      residual $V-\text{target}$ so the square system stays well conditioned
      while Newton polishes the feasible interior.
\item \textbf{Rejection} (\code{MU\_ACCEPT}): a pointwise solution whose
      implied $\mu$ exceeds a physical threshold is rejected and the prior
      value kept, so a spurious ``maximal-binding'' trial cannot corrupt the
      fit. (The threshold is $0.10$ in the no-default model where $\mu$ is
      truly tiny, raised to $0.85$ in the default regime where a $55\%$
      haircut makes the constraint bind hard for real.)
\item \textbf{Rotation} (\S\ref{sec:rotated}): reshaping the box to the
      ergodic cloud removes most infeasible corners at the source.
\end{itemize}
The honest limitation, documented in the replication notes, is that these
mitigate but do not eliminate the problem: a residual $\sim40$--$60\%$ of
axis-aligned corners must be anchored, and their frozen values floor the
achievable residual of the \emph{global} fit. This is intrinsic to a
\emph{global} polynomial basis on a wide box---corner errors are non-local---
and is the reason the notes flag a local/finite-element basis as the
principled next step.

\subsection{Interaction with the outer fixed point}

Finally, the approximation error interacts with the solution algorithm.
Time iteration is successive approximation on the policy; its contraction
rate equals the model's persistence ($\sim0.98$), so it crawls and can
\emph{oscillate}---reaching a good near-fixed-point and then drifting away,
because the two-regime near-unit-root iteration is not globally contractive.
The code's defenses are outer \emph{damping} (update
$V\leftarrow(1-\text{damp})V+\text{damp}\,V^{\text{new}}$), convergence
measured on rule \emph{values} rather than coefficients (the kink makes
coefficients cycle even when values settle), and \emph{best-iterate
tracking} in \code{time\_iterate\_full} (keep the iterate with the fewest
infeasible points and smallest update, not the last). The alternative
Newton solver on the stacked collocation residual
(\code{newton\_collocation.py}) trades the eigenvalue-limited linear rate
for quadratic convergence, at the cost of the kink-handling above.


\section{The branch computation}
\label{sec:branch}

We now come to the piece the reader specifically asked about: how the code
computes ``the branch.'' The term refers to the \emph{two-branch quadrature}
that evaluates the conditional expectations inside the agents' Euler
equations, splitting next period into a no-default branch and a default
branch. Conceptually this is where the priced sovereign risk enters the
model, and it is the reason a \emph{global} solution was built in the first
place.

\subsection{The economic object}

An Euler equation is an equilibrium condition of the form
``price today $=$ expected discounted payoff tomorrow.'' The expectation is
over next period's shocks \emph{and} over whether the sovereign defaults.
Formally, for any payoff $X'$ the banker's/household's condition needs
\begin{equation}
\E_t\big[X'\big] \;=\;
\int\!\!\int\!\!\int
\Big[(1-p^d)\,X'^{\,(d'=0)} + p^d\,X'^{\,(d'=1)}\Big]\;
d\Phi(\epsilon_z)\,d\Phi(\epsilon_g)\,d\Phi(\epsilon_s),
\label{eq:expect}
\end{equation}
with two ingredients to approximate: the integral over the three Gaussian
innovations $(\epsilon_z,\epsilon_g,\epsilon_s)$, and the discrete default
lottery $d'\in\{0,1\}$. Bocola's key modeling point---and the reason the
representative-branch shortcut in the sister \code{code/global/} model
failed---is that this is a genuine \emph{distribution over a discrete future
event}, not a point mass. Sovereign risk raises the \emph{probability weight}
$p^d$ on a low-payoff branch; it does not merely shift a mean.

\subsection{Gauss--Hermite quadrature over the innovations}

The three innovations are independent standard normals, so the integral
over them is done by \emph{Gauss--Hermite} quadrature---the exact tool for
Gaussian weights. An $n$-node Gauss--Hermite rule integrates polynomials up
to degree $2n-1$ exactly against the normal density; the code uses the
probabilists' variant \code{hermegauss} and tensorizes it over the three
shocks:
\begin{verbatim}
def gauss_hermite_3d(n=3):
    x, w = np.polynomial.hermite_e.hermegauss(n)   # E[f], eps ~ N(0,1)
    w = w / w.sum()
    xz, xg, xs = np.meshgrid(x, x, x, indexing="ij")
    wz, wg, ws = np.meshgrid(w, w, w, indexing="ij")
    return xz.ravel(), xg.ravel(), xs.ravel(), (wz*wg*ws).ravel()
\end{verbatim}
With $n=3$ nodes per shock this is $3^3=27$ quadrature nodes---and, note,
this is a \emph{tensor} rule with strictly positive weights, deliberately
kept separate from the Smolyak grid (\S\ref{sec:risks}). Three nodes suffice
because the integrand, once the policy rules are smooth, is well
approximated by a low-degree polynomial in the shocks; the AR(1) laws of
motion map each quadrature node to a next-period exogenous state:
\begin{equation}
\Delta z' = (1-\rho_z)\gamma+\rho_z\Delta z+\sigma_z\epsilon_z,\quad
s' = (1-\rho_s)s^\star+\rho_s s+\sigma_s\epsilon_s, \ \dots
\label{eq:ar1}
\end{equation}
(\code{next\_exog\_full}).

\subsection{Splitting into branches}
\label{sec:branch-split}

The default lottery is the discrete part of \eqref{eq:expect}. The priced
probability of default next period is a \emph{logistic} function of the risk
factor (Bocola's equation 11):
\begin{equation}
p^d(s)\;=\;\frac{1}{1+e^{-s}} .
\label{eq:pd}
\end{equation}
The code computes the future integrand \emph{once per branch} and then takes
the probability-weighted combination. The heart of it is
\code{expect\_pieces\_full} in \code{expectations\_full.py}:
\begin{verbatim}
pd = default_prob(s)                          # weight on the default branch
inv0, lh0, rk0, rb0 = _branch_pieces(S_next, 0, grid, coef, cal)  # d'=0
inv1, lh1, rk1, rb1 = _branch_pieces(S_next, 1, grid, coef, cal)  # d'=1
w0, w1 = (1.0 - pd) * wq, pd * wq             # branch x quadrature weights
return (np.dot(w0, inv0) + np.dot(w1, inv1),  # E[1/C']
        np.dot(w0, lh0)  + np.dot(w1, lh1),   # E[w'/C']
        np.dot(w0, rk0)  + np.dot(w1, rk1),   # capital-return integrand
        np.dot(w0, rb0)  + np.dot(w1, rb1))   # bond-return integrand
\end{verbatim}
Read this carefully, because it is the whole mechanism. \code{w0} and
\code{w1} are outer products of the Gauss--Hermite weights \code{wq} with the
scalar branch probabilities $1-p^d$ and $p^d$: each of the $27$ shock nodes
is duplicated into a no-default copy and a default copy, and the two weight
vectors sum to the full expectation. The four returned numbers are exactly
the conditional-expectation ``pieces'' that the Euler residuals need---
$\E[1/C']$, $\E[w'/C']$, and the capital- and bond-return integrands---
evaluated over the joint distribution of shocks \emph{and} default.

\subsection{What differs across branches}

The two branches share the same next-period \emph{shock} nodes but evaluate
different objects. \code{\_branch\_pieces} makes the distinction precise:
\begin{verbatim}
def _branch_pieces(S_next, d_next, grid, coef, cal):
    S_eval = grid.clip(S_next)                       # never extrapolate
    Cn  = np.maximum(grid.eval(coef["C"][d_next],  S_eval), 1e-8)
    QBn = np.maximum(grid.eval(coef["QB"][d_next], S_eval), 1e-4)
    ...
    surv  = 1.0 - d_next * cal["haircut"]            # 1 if d'=0, 1-D if d'=1
    RBpay = surv * (pi + (1-pi)*(iota + QBn))        # haircut bond payoff
    ...
\end{verbatim}
Three things branch:
\begin{enumerate}
\item \textbf{The rules used.} Decision rules carry a \emph{separate
      coefficient vector per regime}, \code{coef[rule] = (coef\_d0,
      coef\_d1)}. Next period's policy is read from whichever regime $d'$
      realizes---so the default branch continuation uses the crisis-regime
      rules, the no-default branch the normal-regime rules.
\item \textbf{The bond payoff.} In the default branch the survival share is
      $\text{surv}=1-D$ with $D=0.55$ (the Greek PSI haircut), so the bond
      pays only $45\%$ of its promise. This haircut, weighted by $p^d(s)$,
      is the direct channel by which higher priced risk lowers today's bond
      price $Q_B$.
\item \textbf{The current-period statics.} The \emph{current} regime
      $d_{\text{cur}}$ (distinct from next period's $d'$) enters the current
      period map through the same $\text{surv}$ factor---a haircut on bonds
      held \emph{into} a default period.
\end{enumerate}
Re-default is allowed: $d'$ is drawn independently of the current regime, so
both current regimes share the same two-branch continuation. This is the
``distribution, not a point mass'' structure spelled out in the file header.

\subsection{From branch pieces to prices and the multiplier}

The four expectation pieces close the model. Back in
\code{point\_block\_full}, they are scaled by the discount factor
$\beta\,C\,e^{-\Delta z}$ into $\E[\Lambda']$, $\E[\hat\Lambda']$ and the two
return integrands, and combined into three unit-free Euler residuals plus
the closed-form multiplier
\begin{equation}
\mu \;=\; \max\!\Big\{\,1-\frac{\E[\hat\Lambda']\,R\,N}{\lambda(Q_KK'+Q_BB')}\,,\;0\Big\}
\qquad(\code{mu\_closed\_form}),
\label{eq:mu}
\end{equation}
which is Bocola's Proposition-1 formula. The economic pass-through lives
entirely in \eqref{eq:pd}--\eqref{eq:mu}: a rise in the risk factor $s$
raises $p^d$, which raises the weight \code{w1} on the haircut branch, which
lowers the bond-return integrand \code{rb}, which lowers the equilibrium
bond price $Q_B$ that clears the bond Euler; the resulting mark-to-market
loss on bank balance sheets lowers net worth $N$, which through \eqref{eq:mu}
tightens the constraint ($\mu\uparrow$), which widens lending spreads. The
replication reproduces exactly this signature: at Bocola's $2.5\%$ default
probability the bond price falls $-12\%$ and net worth falls
(Figure~\ref{fig:branch}), the correct contractionary sign.

\begin{figure}[h]
\centering
\includegraphics[width=\textwidth]{branch.pdf}
\caption{\emph{Left:} the two-branch quadrature. The current state is pushed
through the AR(1) laws to $27$ Gauss--Hermite shock nodes, each of which
splits into a no-default continuation (weight $1-p^d$) and a haircut default
continuation (weight $p^d$); the expectation is the probability- and
quadrature-weighted sum. \emph{Right:} the priced default probability
$p^d(s)=1/(1+e^{-s})$ that sets the branch weights, from its
balanced-growth value $s^\star=-7.06$ up to the stress region.}
\label{fig:branch}
\end{figure}

\subsection{Why the branch structure is the whole point}

It is worth stating plainly why this machinery exists. The sister
two-country model in \code{code/global/} approximates the same risk with a
single \emph{representative} post-default branch computed once and reused---a
point-mass shortcut. That approximation gets the sign \emph{wrong}: it makes
the crisis-state bank lever \emph{into} capital, so sovereign risk comes out
expansionary. Bocola's model has no separate ``branch'' to bolt on: the bad
state is \emph{the same decision rules evaluated at a different point of the
state space}, reached with probability $p^d(s)$. Representing that
faithfully requires (i) a global solution defined over the whole state
space---hence Chebyshev--Smolyak---and (ii) the genuine two-branch
quadrature above. The two halves of this note are therefore not independent:
the sparse-grid interpolation exists precisely so that
\eqref{eq:expect}--\eqref{eq:pd} can be evaluated at arbitrary future
states, and the branch quadrature is what turns those interpolated rules
into the priced-risk pass-through.


\section{Assembling the solver}
\label{sec:assemble}

To place the pieces in the full loop, the solve is a fixed point over the
rule coefficients:
\begin{enumerate}
\item \textbf{Seed.} Start from a first-order perturbation solution (a smooth,
      globally-roughly-correct guess), fitted onto the Smolyak grid. A
      constant-guess seed instead leaves a third of points infeasible.
\item \textbf{Inner solve.} At each of the $85$ (or $389$) grid points, and
      each regime $d\in\{0,1\}$, solve the $3$-equation Euler system for
      $(C,R,Q_B)$ by bounded least squares, using the two-branch quadrature
      of \S\ref{sec:branch} to form the expectations against the
      \emph{current} coefficient vectors.
\item \textbf{Fit.} Convert the updated point values back to Chebyshev
      coefficients with the cached LU factor (\code{grid.fit}).
\item \textbf{Damp and iterate.} Damp the update, measure convergence on
      rule values, track the best iterate, and repeat until the update norm
      falls below tolerance. Optionally polish with the FB-smoothed Newton
      solver.
\item \textbf{Validate.} Simulate the solved rules forward to trace the
      ergodic set, and report the Euler residual \emph{along the path the
      model actually visits}. The accepted accuracy standard---Bocola's
      own---is accuracy on the ergodic set, which the package meets (mean
      $\log_{10}$ residual $\approx-3.1$, i.e.\ $\sim0.1\%$), even though a
      wide-box global fit cannot be uniformly machine-accurate in the
      never-visited infeasible corners.
\end{enumerate}


\section{Summary}
\label{sec:summary}

The computation rests on three theorems and one modeling insight.
\emph{Chebyshev interpolation} on Lobatto nodes is near-minimax and
spectrally accurate for smooth functions---so a policy function becomes a
short, stable coefficient vector---while being explosively wrong under
extrapolation and only algebraically accurate at kinks. \emph{Smolyak sparse
grids} lift that construction to six dimensions at polynomial rather than
geometric cost, by keeping only low-total-level tensor blocks of nested
Chebyshev levels, at the price of a faster-growing Lebesgue constant and
exactness on only the hyperbolic-cross polynomial space. \emph{Gauss--Hermite
quadrature} evaluates the Gaussian expectations exactly to high polynomial
degree with positive weights. And the \emph{two-branch structure} is the
economic heart: sovereign risk is a probability weight $p^d(s)$ on a
haircut-default continuation, evaluated by interpolating the crisis-regime
rules at stochastically generated future states. The recurring discipline
that makes all of this work is respect for the region the model actually
visits---clip before you extrapolate, rotate the box onto the ergodic cloud,
smooth every kink the Jacobian sees, anchor the infeasible corners, and
judge accuracy on the ergodic set. Each safeguard in the code maps to a
specific failure mode in the theory above.


\section*{Where to read the code}
\addcontentsline{toc}{section}{Where to read the code}
\begin{center}
\small
\begin{tabular}{@{}ll@{}}
\toprule
Concept & File / function \\
\midrule
Chebyshev recurrence, nodes & \code{smolyak.py}: \code{chebyshev\_basis\_1d}, \code{\_level\_points} \\
Nested levels, new points/degrees & \code{smolyak.py}: \code{\_new\_points}, \code{\_new\_degrees} \\
Smolyak combination rule & \code{smolyak.py}: \code{\_multi\_indices}, \code{SmolyakGrid.\_\_init\_\_} \\
Fit / eval (LU-cached collocation) & \code{smolyak.py}: \code{SmolyakGrid.fit}, \code{.eval} \\
Box mapping, clipping & \code{smolyak.py}: \code{to\_unit}/\code{from\_unit}, \code{clip} \\
PCA-rotated grid & \code{rotated\_grid.py}: \code{RotatedGrid} \\
Gauss--Hermite quadrature & \code{expectations\_full.py}: \code{gauss\_hermite\_3d} \\
Two-branch expectation & \code{expectations\_full.py}: \code{expect\_pieces\_full}, \code{\_branch\_pieces} \\
Logistic default probability & \code{expectations\_full.py}: \code{default\_prob} \\
Euler residuals, $\mu$ closed form & \code{time\_iteration\_full.py}, \code{period\_map.py}: \code{mu\_closed\_form} \\
Fischer--Burmeister smoothing & \code{newton\_collocation.py}: \code{\_fb}, \code{\_point\_residual\_fb} \\
Outer fixed point, best iterate & \code{time\_iteration\_full.py}: \code{time\_iterate\_full} \\
Ergodic simulation / box adapt & \code{simulate.py}, \code{simulate\_full.py} \\
\bottomrule
\end{tabular}
\end{center}


\section*{References and further reading}
\addcontentsline{toc}{section}{References and further reading}
\small
\begin{itemize}
\item L.\ Bocola (2016), ``The Pass-Through of Sovereign Risk,''
      \emph{Journal of Political Economy} 124(4). The model and its
      global-solution method (\S II.C, eqs.\ 11--12 for the default
      process).
\item S.\ A.\ Smolyak (1963), ``Quadrature and interpolation formulas for
      tensor products of certain classes of functions,'' \emph{Soviet Math.\
      Dokl.} 4. The original sparse-grid construction.
\item D.\ Krueger and F.\ Kubler (2004), ``Computing equilibrium in OLG
      models with stochastic production,'' \emph{J.\ Economic Dynamics \&
      Control} 28. Introduced Smolyak sparse grids to economics; the nested
      construction the code follows.
\item K.\ Judd, L.\ Maliar, S.\ Maliar and R.\ Valero (2014), ``Smolyak
      method for solving dynamic economic models: Lagrange interpolation,
      anisotropic grid and adaptive domain,'' \emph{J.\ Economic Dynamics \&
      Control} 44. Anisotropic levels and the efficient implementation.
      \url{https://bfi.uchicago.edu/wp-content/uploads/Judd-Maliar-Valero-1.pdf}
\item L.\ N.\ Trefethen (2013), \emph{Approximation Theory and Approximation
      Practice}, SIAM. The definitive modern treatment of Chebyshev
      interpolation, the Runge phenomenon, Lebesgue constants and spectral
      convergence (Chebfun examples on Lebesgue constants:
      \url{https://www.chebfun.org/examples/approx/LebesgueConst.html}).
\item J.\ Brumm and S.\ Scheidegger (2017), ``Using adaptive sparse grids to
      solve high-dimensional dynamic models,'' \emph{Econometrica} 85.
      Adaptive/local sparse grids---the principled route past the
      global-basis corner problem of \S\ref{sec:risks}.
\item Package internal notes: \code{docs/bocola2016\_replication.md} (BGP
      derivation, stage-by-stage accuracy, the corner/convergence
      diagnostics summarized here).
\end{itemize}

\end{document}
